\section{Concrete Configurations}
\label{sec:configurations}

This section catalogues seven concrete configurations as labeled
points in the taxonomy of \S\ref{sec:dimensions}. The labels C1
through C7 used in the comparison tables of \S\ref{sec:analysis} are
fixed here. Each subsection states the six axis values, the
relation(s) the configuration commonly carries, and a short fit
summary that cross-references the relevant cells of
Tables~\ref{tab:cost}--\ref{tab:pq}. Section~\ref{sec:migration}
discusses migration paths between configurations;
Section~\ref{sec:recommendations} maps deployment profiles to them.

Six of the seven (C1--C6) are deployable today on the chain named in
their first axis. C7 is included as an undeployable point: the
configuration is cryptographically defined and supported by mature
proving libraries, but is blocked at the host-function layer of its
named anchor. We retain C7 in the catalogue rather than relegate it
to the §6 migration discussion alone because it is the
shortest-distance witness to the host-function axis being a
binding constraint independent of the other five axes; without C7,
the §3 dependency analysis would lack a direct visible counterexample
to the assumption that a transparent-setup PQ-ready SNARK can be
deployed wherever the underlying construction is mature.

\subsection{C1 --- Stellar + BLS12-381 + PLONK + Poseidon}
\label{sec:cfg-c1}

\textbf{Axes.} Anchor chain: Stellar (Soroban). SNARK: PLONK with
KZG commitments~\cite{gabizon-2019-plonk}. Curve / in-circuit hash:
BLS12-381 / Poseidon, with SHA-256 as the outer commitment binding.
PQ stance: pairing-based, broken under Shor. Trusted setup:
\emph{universal updatable, instantiated by reuse of the Ethereum
Foundation KZG Summoning Ceremony}~\cite{eth-kzg-ceremony}.
Verification host: native BLS12-381 pairing, $\mathbb{G}_1$
addition, and $\mathbb{G}_1$ scalar-multiplication host functions
added in Stellar Protocol 22~\cite{stellar-cap-0059}, with
Poseidon / Poseidon2 permutation primitives over BLS12-381's scalar
field added in Stellar Protocol 25~\cite{stellar-cap-0075}.

\textbf{Relations.} The SEP draft at~\cite{sep-onym} carries a
membership relation $R_{\mathsf{Mem}}$ and an update relation
$R_{\mathsf{Upd}}$ binding both the old and new commitments and the
epoch counter; per-policy variants (single-admin, single-member,
quorum, fixed-set) are realized as distinct policy-instances
deployed alongside the membership relation. The reference C1
deployment splits the five policy-instances across five separate
Soroban contracts (one policy per contract); each contract holds
two relations (membership + update) compiled at three circuit-size
tiers (small, medium, large, indexed by anticipated membership-set
cardinality), giving a total of thirty verifying keys across the
deployment. All thirty verifying keys share a single SRS --- this
is the headline property of the universal-updatable family that
distinguishes C1's trust posture from Groth16-shaped configurations
such as C2. The §4 benchmark harness treats each contract as a
distinct measurement unit and aggregates per-contract numbers into
the single C1 row of the comparison tables. The five-contract
split is a deployment choice, not a property of the registry
abstraction; an alternative C1 implementation could dispatch on a
policy tag from a single contract.

\textbf{Trust supply chain for the SRS.} The C1 deployment pins the
provenance of its trusted-setup parameters as an explicit
hash-chained sequence: the upstream ceremony transcript SHA-256 is
recorded at build time; an extracted binary form of the SRS is
hashed; per-tier verifying keys baked from the SRS are hashed and
mirrored into a cross-platform test-vector file; and each contract's
deployed WASM bytecode includes the per-tier verifying key bytes
directly, so the contract's WASM-hash transitively binds the
verifying key. A third party can re-fetch the EF transcript,
re-extract, re-bake, and check every hash in the chain against the
deployed contract's WASM. The reproducibility property the chain
provides --- that the deployed contracts use the canonical EF KZG
SRS and not a substituted local one --- is the operational
counterpart to the §3.5 ``one honest contributor suffices'' formal
property: the formal property is unconditional in the model, the
supply-chain pinning is what gives a reviewer a deterministic
artifact to verify it against.

\textbf{Fit summary.} Cheap-end of the verification-cost regime
(Table~\ref{tab:cost}), with the §4.1 caveat that PLONK's
multi-scalar-multiplications add overhead beyond Groth16's
single-pairing-check verifier shape; on Stellar this overhead is
attenuated by CAP-0059's $\mathbb{G}_1$ scalar-mul host functions
but is not eliminated. Proof size ${\sim}700$~B
(Table~\ref{tab:proof-size}), an order of magnitude larger than
Groth16's 192~B but two orders smaller than the FRI-based
configurations. The trust posture beyond the SRS is the
SNARK's algebraic-group-model heuristics; the SRS's residual trust
is one-honest-of-${\sim}141{,}000$ contributors~\cite{eth-kzg-ceremony},
which is the universal-updatable family's structural advantage and
which has no analogue in the Groth16-with-per-circuit-MPC
configurations. Not post-quantum: the KZG commitment's binding
property reduces to discrete log on $\mathbb{G}_1$.

\subsection{C2 --- Ethereum L1 + BN254 + Groth16 + Poseidon}
\label{sec:cfg-c2}

\textbf{Axes.} Anchor chain: Ethereum L1. SNARK: Groth16.
Curve / hash: BN254 / Poseidon (Pedersen-or-MiMC in older
deployments). PQ stance: pairing-based. Trusted setup: per-circuit
MPC. Verification host: BN254 pairing precompile (EIP-197) and G1
add / scalar-mul precompiles (EIP-196).

\textbf{Relations.} The original Tornado Cash deposit-spend
construction~\cite{pertsev-2019-tornado} fits this configuration with
a single-relation circuit binding the deposit-tree root, a nullifier,
and an external recipient. Semaphore~\cite{semaphore-2024} fits it
with an identity-tree relation parameterized by an external nullifier
for one-shot signaling.

\textbf{Fit summary.} The classical configuration in the literature.
Verification cost dominated by the BN254 pairing precompile and one
G1 MSM; well-characterised at $\sim$200K--250K gas for a single
proof. The
Kim--Barbulescu attack~\cite{kim-barbulescu-2016} reduced the
practical security of BN254 from the originally targeted $\sim$128
classical bits to $\sim$100 bits, which is the security-relevant
caveat on this configuration's continued deployment, independent of
its post-quantum posture. EIP-2537 (BLS12-381 precompiles on Ethereum L1) would
unlock a migration to the same security level as C1 but at the time
of writing has not been activated.

\subsection{C3 --- Aztec + UltraPlonk + BN254}
\label{sec:cfg-c3}

\textbf{Axes.} Anchor chain: Aztec, an L2 SNARK-rollup whose
execution model is itself a SNARK. SNARK: UltraPlonk (PLONK with
custom gates and lookup arguments) implemented in Noir. Curve / hash:
BN254 / Pedersen-or-Poseidon. PQ stance: pairing-based. Trusted
setup: universal updatable (Aztec runs an Aztec-specific Phase-1
ceremony that all relations share). Verification host:
verifier-as-rollup --- the chain itself is the verifier.

\textbf{Relations.} Application-defined; Aztec contracts compile to
private functions whose witnesses include the relation and whose
output is a proof checked by the rollup as a precondition for state
transition.

\textbf{Fit summary.} The verifier-as-rollup posture means that
application-level proofs ride the same machinery that validates the
chain, so verification cost is amortized across the rollup's batched
proof aggregation. This is qualitatively different from C1 / C2
verification cost and is reported in
Table~\ref{tab:cost} in L2 gas units rather than as a single
on-chain fee. Trade-offs: the sequencer is centralized in the
current alpha, which dominates the trust posture; proofs are larger
than Groth16 but smaller than FRI; and the entire deployment is a
single-vendor stack at the time of writing.

\subsection{C4 --- Mina + Pickles + Pasta}
\label{sec:cfg-c4}

\textbf{Axes.} Anchor chain: Mina. SNARK: Pickles, Mina's recursive
proof-composition system based on Halo-style accumulator
techniques~\cite{bowe-2019-halo}. Curve / hash:
Pasta (Pallas/Vesta cycle)~\cite{hopwood-2020-pasta} / Poseidon. PQ
stance: classical-EC non-pairing, broken under Shor.
Trusted setup: transparent (Pickles uses no setup). Verification host:
Mina's protocol-level Pickles verifier, equivalent to a built-in
host function.

\textbf{Relations.} Mina applications compose recursive proofs;
zkApps express their state-transition predicate in OCaml, compiled to
Pickles circuits.

\textbf{Fit summary.} The transparent-setup-with-built-in-host
combination is rare in the surveyed configurations. Recursion is
structural: every Mina block's state transition is itself a Pickles
proof, so application-level proofs piggyback on the chain's recursive
verification. Proofs are 10--20~KB owing to recursion overhead.
Mina's $\sim$300-block-producer validator set sits between Stellar's
$\sim$30 and Ethereum's $\sim$10$^6$ on the decentralization axis.

\subsection{C5 --- L2 + Plonky2/3 + KZG over BN254}
\label{sec:cfg-c5}

\textbf{Axes.} Anchor chain: an Ethereum L2 (Polygon zkEVM in the
canonical instance, but the configuration is L2-generic). SNARK:
Plonky2 / Plonky3~\cite{plonky2-2022,plonky3-spec} for the
inner-prover-time-optimized layer, with a final KZG wrap over BN254
for L1 settlement compatibility. Curve / hash: Goldilocks at the
inner layer, BN254 at the outer wrap; Poseidon is the canonical
in-circuit hash. PQ stance: \emph{pairing-based} as deployed,
because the outer KZG wrap is what L1 actually verifies and
therefore dominates the configuration's post-quantum class; the
inner FRI prover is post-quantum-secure as a sub-component but is
not the bottleneck (\S\ref{sec:analysis-pq}). Trusted setup:
universal updatable (KZG ceremony shared with broader BN254
ecosystem). Verification host: the L2 sequencer's batched proof
verification on L1, ultimately calling the BN254 pairing precompile.

\textbf{Relations.} Application-defined; the L2 architecture
typically allows applications to ride the L2's general-purpose ZK
execution.

\textbf{Fit summary.} Frequently mis-classified as post-quantum
because the inner prover is FRI-based; the post-quantum table
(Table~\ref{tab:pq}) places it at level zero because the outer wrap
is the verification-relevant component. Wins on prover wall-clock
(Table~\ref{tab:prover-time}) by virtue of the FRI inner layer.

\subsection{C6 --- STARK over Goldilocks (anchor-agnostic)}
\label{sec:cfg-c6}

\textbf{Axes.} Anchor chain: variable; canonical instances anchor on
StarkNet (an L2 whose verifier is FRI-friendly), but the same
configuration could anchor on a chain that subsidizes large-proof
verification. SNARK: STARK~\cite{ben-sasson-2018-stark}.
Curve / hash: Goldilocks (no curve) / a STARK-friendly hash family
(Rescue-Prime / Poseidon2 with appropriate rate). PQ stance:
hash-based, conjectured PQ-secure under standard collision-resistance
assumptions. Trusted setup: transparent. Verification host:
chain-dependent --- StarkNet provides a built-in verifier; on a
non-FRI-aware chain, verification is pure bytecode and prohibitively
expensive.

\textbf{Relations.} Application-defined; STARK relations are
typically expressed in Cairo or a similar STARK-friendly DSL.

\textbf{Fit summary.} The unique configuration in the surveyed set
that combines transparent setup, hash-based post-quantum security,
and a deployable verification host. Proof sizes 50--200~KB
(Table~\ref{tab:proof-size}) are the dominant cost; verification
cost is moderate on FRI-aware anchors and prohibitive elsewhere.

\subsection{C7 --- Stellar + Plonky3 + FRI (hypothetical, blocked)}
\label{sec:cfg-c7}

\textbf{Axes.} Anchor chain: Stellar. SNARK: Plonky3 / FRI. Curve / hash:
Goldilocks / Poseidon2 (no curve at the verification layer). PQ
stance: hash-based, conjectured PQ-secure. Trusted setup: transparent.
Verification host: \textbf{none} --- Stellar Protocol 22 added
BLS12-381 host functions but no FRI host functions, no
NTT host function, and no Goldilocks-arithmetic host function. A
verifier for this configuration would run in pure WASM, with FRI
folding and Merkle-path verification interpreted at WASM-level cost.

\textbf{Relations.} Identical in shape to C1 (a membership and update
pair) but realized in a Plonky3-style transparent-setup circuit.

\textbf{Fit summary.} Included in the catalogue not as a deployment
target but as the visible point that makes the host-function axis
binding. Pure-WASM verification of FRI on Stellar is estimated
at one to two orders of magnitude more expensive than C1 in
verification gas (Table~\ref{tab:cost}, model entry); in proof size
it matches C6. The configuration becomes deployable only under the
host-function gap-closing scenarios in
\S\ref{sec:migration}.

\subsection{Configurations not in the catalogue}

The seven configurations cover the deployment-relevant points in the
six-axis taxonomy of \S\ref{sec:dimensions} that we found documented
in the open literature or in production-grade open-source
implementations. We deliberately omit several adjacent points:

\begin{itemize}
  \item \textbf{Bitcoin + inscription-style data attachment}: BitVM
    and adjacent schemes can in principle host SNARK verifiers, but
    the verifier-host axis is sufficiently degenerate (verification
    requires multi-step interactive challenges) that we treat it as a
    distinct system class outside the registry abstraction.
  \item \textbf{Solana + Light Protocol}: a real deployment that
    fits the abstraction, omitted for space. Its SNARK-system axis
    matches C2 (Groth16-shape), but anchor chain, host-function
    profile (Solana lacks BN254 / BLS12-381 precompiles at the
    eBPF/SBF level and verifies via per-instance compiled programs),
    fee market, and validator-set composition all diverge from C2.
    A surveyed treatment is left to future work.
  \item \textbf{Lattice-based SNARKs}: no production deployment we
    are aware of; flagged as an open problem in \S\ref{sec:open}.
  \item \textbf{No-chain (DA-layer-only)}: the
    advancement abstraction collapses without a single-writer
    consensus to enforce monotonicity; we treat this as a different
    primitive.
\end{itemize}
